% ===============================================================
%  Chapter 3 -- Parallel Processing and Multicore Systems   (Part I)
% ===============================================================
\chapter{Parallel Processing and Multicore Systems}
\label{ch:parallel}

\begin{figure}[htbp]
\centering
\adjustbox{max width=\linewidth}{%
\begin{forest}
  chapter map,
  [Chapter 3
    [From one core to many [Frequency wall] [MIMD multicores]]
    [Shared memory {\&} threads [Fork--join in OpenMP]]
    [Interconnects [{Bus, ring, mesh}]]
    [Amdahl's law [Speedup ceiling] [Why coherence is needed]]
  ]
\end{forest}}
\caption[]{Structure of Chapter~3: why chips became
multicore, how they are programmed and connected, what limits the gain from more
cores, and where coherence becomes necessary.}
\label{fig:ch3-map}
\end{figure}

% ---------------------------------------------------------------
\section{From One Core to Many}
\label{sec:par-why}
\label{sec:par-defs}
\label{sec:par-flynn}
\sectionmap{[One core is not enough [Clock speed stalled near 3 GHz]
  [Little more parallelism in one program] [Spend transistors on more cores]]}
\index{Sarangi textbook}%

Until the mid-2000s an unchanged program ran faster on each new processor. Then
clock frequency levelled off at around 3\,GHz, because beyond that point power and
heat grew faster than performance, and wide superscalar cores could no longer find
enough independent instructions to issue~\cite[\S12.1]{sarangi2025basic}. Moore's
law~\cite{moore1965cramming} still gave designers more transistors each
generation, so they spent them on several simpler cores in place of one bigger
core.

Sarangi defines \emph{multiprocessing} as several processors working in parallel,
and a \emph{multicore processor}\index{multicore processor} as a multiprocessor
whose processors (cores) share one chip~\cite[\S12.1, Def.~116]{sarangi2025basic}.
In Flynn's taxonomy~\cite{flynn1972organizations}\cite[\S12.3]{sarangi2025basic} a
multicore is a MIMD machine, since every core fetches its own instructions and
works on its own operands. Multiprocessors came first. By the late 1990s
motherboards carried four or eight processors, IBM's POWER4 put two cores on one
chip in 2001, and Intel and AMD followed in 2005~\cite[\S12.1]{sarangi2025basic}.
This report deals with the multicore case, where the cores share one chip and one
physical address space.

% ---------------------------------------------------------------
\section{Shared Memory, Threads and Fork--Join}
\label{sec:par-sm-mp}
\label{sec:par-threads}
\sectionmap{[Programming model [Shared memory: common load and store]
  [Message passing: explicit send] [Fork--join threads share data]]}
\index{Sarangi textbook}%

Parallel programs communicate in one of two ways (\cref{fig:sm-mp}). Threads in a
shared-memory program use one address space and exchange data with ordinary loads
and stores. Processes in a message-passing program, such as one written with MPI,
have private memories and exchange data only through explicit \texttt{send} and
\texttt{receive} calls~\cite[\S12.2.2]{sarangi2025basic}. Coherence is needed only
in the shared-memory case, because only there can two cores cache the same
location.

\begin{figure}[htbp]
\centering
\begin{bigfig}[0.85]
\begin{tikzpicture}[font=\sffamily\scriptsize,
    u/.style={draw, thick, rounded corners=2pt, minimum width=1.3cm, minimum height=0.55cm},
    m/.style={draw, thick, rounded corners=2pt, fill=leafC!14, draw=leafC,
              minimum height=0.6cm, align=center}]
  \node[font=\sffamily\bfseries\small, text=tealC!80!black] at (2.0,2.7) {Shared memory};
  \foreach \i/\c in {0/tealC,1/plumC,2/amberC} { \node[u, fill=\c!15, draw=\c] (t\i) at (\i*2.0,1.9) {thread \i}; \draw[thick, -{Stealth}, \c] (t\i) -- (\i*2.0,0.95); }
  \node[m, minimum width=5.4cm] at (2.0,0.6) {one shared address space};
  \node[text=coreGrey] at (2.0,-0.2) {communicate by \texttt{load} / \texttt{store}};
  \begin{scope}[xshift=7.4cm]
  \node[font=\sffamily\bfseries\small, text=brickC] at (2.4,2.7) {Message passing};
  \foreach \i in {0,1,2} {
    \node[u, fill=slateC!14, draw=slateC] (p\i) at (\i*2.4,1.9) {process \i};
    \node[m, minimum width=1.3cm] (mm\i) at (\i*2.4,0.6) {private\\memory};
    \draw[thick] (p\i) -- (mm\i);
  }
  \draw[-{Stealth}, thick, brickC] (p0) -- node[above=2pt, font=\sffamily\tiny]{\texttt{send}} (p1);
  \draw[-{Stealth}, thick, brickC] (p1) -- node[above=2pt, font=\sffamily\tiny]{\texttt{send}} (p2);
  \node[text=coreGrey] at (2.4,-0.2) {communicate only by messages};
  \end{scope}
\end{tikzpicture}
\end{bigfig}
\caption[Shared memory and message passing]{The two programming models. Only the
left-hand model has locations that two cores can both cache.}
\label{fig:sm-mp}
\end{figure}

Sarangi defines a thread as a sub-program with its own program counter and stack
that runs in the same address space as other threads of the
program~\cite[\S12.2.2, Def.~120]{sarangi2025basic}. Shared-memory programs usually
follow the fork--join pattern. A sequential section starts $N$ threads, each thread
works on its slice of the data, the program waits for all of them, and then it
combines their results. \Cref{lst:omp} shows this in \gls{openmp}.

\begin{lstlisting}[language=C, caption={Fork--join in OpenMP: summing an array in parallel.}, label=lst:omp]
double partial[NTHREADS];               /* shared array of results  */
#pragma omp parallel num_threads(NTHREADS)
{
    int me = omp_get_thread_num();      /* which slice is mine?     */
    int first = me * (N / NTHREADS), last = first + N / NTHREADS;
    double s = 0;
    for (int k = first; k < last; k++) s += a[k];
    partial[me] = s;                    /* write a shared location  */
}                                       /* implicit join here       */
double sum = 0;
for (int t = 0; t < NTHREADS; t++) sum += partial[t];
\end{lstlisting}

The statement \texttt{partial[me] = s} writes a shared array that another core
reads after the join. The program says nothing about how the value travels to the
reading core, so the hardware has to handle it.

% ---------------------------------------------------------------
\section{Interconnection Networks}
\label{sec:par-noc}
\sectionmap{[Interconnects [Bus: one shared medium] [Ring and mesh: point-to-point]
  [Diameter and bisection bandwidth]]}
\index{network on chip}%

Cores and caches exchange messages over an on-chip interconnect, or
\gls{noc}\index{network on chip}. Two numbers describe it. The diameter is the
longest shortest path between two nodes, which sets the worst-case latency. The
bisection bandwidth is the smallest number of links that must be cut to split the
network in half, which limits throughput~\cite[\S12.6.2]{sarangi2025basic}.
\Cref{fig:topologies,tab:topologies} compare the topologies used in multicores.

\begin{figure}[htbp]
\centering
\begin{minipage}[c]{0.66\linewidth}\centering
\adjustbox{max width=\linewidth}{%
\begin{tikzpicture}[font=\sffamily\scriptsize,
    n/.style={circle, thick, minimum size=0.42cm, inner sep=0pt}]
  \node[font=\sffamily\bfseries\small, text=brickC] at (1.5,2.9) {Bus};
  \foreach \i in {0,...,3} { \node[n, draw=brickC, fill=brickC!18] (b\i) at (\i*1.0,2.0) {\i}; \draw[thick] (b\i) -- (\i*1.0,1.2); }
  \draw[very thick, brickC!70!black] (-0.4,1.2) -- (3.4,1.2);
  \node[text=coreGrey, align=center] at (1.5,0.6) {one message at a time;\\every node sees it};
  \begin{scope}[xshift=5.4cm, yshift=1.5cm]
  \node[font=\sffamily\bfseries\small, text=tealC!80!black] at (0,1.4) {Ring};
  \foreach \i in {0,...,7} { \node[n, draw=tealC, fill=tealC!18] (r\i) at ({90-\i*45}:1.0) {\i}; }
  \foreach \i [evaluate=\i as \j using {int(mod(\i+1,8))}] in {0,...,7} { \draw[thick, tealC!80!black] (r\i) -- (r\j); }
  \end{scope}
  \begin{scope}[xshift=7.6cm, yshift=0.5cm]
  \node[font=\sffamily\bfseries\small, text=plumC] at (1.2,2.4) {2D mesh};
  \foreach \x in {0,...,3} \foreach \y in {0,...,3} {
    \node[n, minimum size=0.32cm, draw=plumC, fill=plumC!18] (m\x\y) at (\x*0.8,\y*0.6) {};
  }
  \foreach \x in {0,...,3} \foreach \y [count=\yy] in {0,...,2} { \draw[thick, plumC!70] (m\x\y) -- (m\x\yy); }
  \foreach \y in {0,...,3} \foreach \x [count=\xx] in {0,...,2} { \draw[thick, plumC!70] (m\x\y) -- (m\xx\y); }
  \end{scope}
\end{tikzpicture}}
\captionof{figure}[Three interconnect topologies]{A bus, an eight-node ring and a
$4\times4$ mesh. A bus delivers each message to every node; rings and meshes
carry point-to-point messages and scale further.}
\label{fig:topologies}
\end{minipage}\hfill
\begin{minipage}[c]{0.31\linewidth}\centering\small
\captionof{table}{Diameter and bisection bandwidth for $N$ nodes~\cite[\S12.6.3]{sarangi2025basic}.}
\label{tab:topologies}
\setlength{\tabcolsep}{3pt}
\begin{tabular}{@{}lcc@{}}
\toprule
 & Diam. & Bisect. \\
\midrule
Ring & $N/2$ & 2 \\
Mesh & $2\sqrt{N}-2$ & $\sqrt{N}$ \\
Torus & $\approx\sqrt{N}$ & $2\sqrt{N}$ \\
Fat tree & $2\log N$ & $N/2$ \\
\bottomrule
\end{tabular}
\end{minipage}
\end{figure}

A bus has a property the other topologies do not get for free. Only one message
is on it at a time, so every node sees every message in the same order.
Chapter~4 uses this property.

% ---------------------------------------------------------------
\section{Amdahl's Law and the Need for Coherence}
\label{sec:par-amdahl}
\label{sec:par-coherence}
\sectionmap{[Amdahl's law [Parallel fraction f] [Speedup saturates at 1/(1-f)]
  [Private caches + shared memory = coherence]]}
\index{Amdahl's law}%

If a fraction $f$ of the work spreads perfectly over $p$ cores and the rest runs on
one core, the speedup is~\cite{amdahl1967validity}\cite[\S12.2.3]{sarangi2025basic}
\begin{equation}
  S(p) = \frac{1}{(1-f) + f/p},
  \label{eq:amdahl}
\end{equation}
which tends to $1/(1-f)$ as $p$ grows (\cref{fig:amdahl}). A system close to this
ceiling can behave quite differently from the same system with two cores, which is
why our experiments sweep 2, 4, 8 and 16 cores.

\begin{figure}[htbp]
\centering
\begin{tikzpicture}
  \begin{axis}[width=10.5cm, height=5cm,
      xlabel={Number of cores, $p$}, ylabel={Speedup $S(p)$},
      xmin=1, xmax=64, ymin=0, ymax=16,
      xtick={1,2,4,8,16,32,64}, xticklabels={1,2,4,8,16,32,64},
      xmode=log, log basis x=2,
      legend pos=north west, legend style={font=\small},
      grid=major, grid style={coreGrey!15},
      label style={font=\small}, tick label style={font=\small}]
    \addplot[very thick, tealC, domain=1:64, samples=100] {1/((1-0.95)+0.95/x)};
    \addlegendentry{$f=0.95$}
    \addplot[very thick, plumC, domain=1:64, samples=100] {1/((1-0.80)+0.80/x)};
    \addlegendentry{$f=0.80$}
    \addplot[very thick, amberC, domain=1:64, samples=100] {1/((1-0.50)+0.50/x)};
    \addlegendentry{$f=0.50$}
  \end{axis}
\end{tikzpicture}
\caption[Amdahl's law]{\Cref{eq:amdahl} for three parallel fractions. With
95\,\% of the program parallel, speedup can never pass $1/(1-f)=20$, and at 16
cores it is well below that.}
\label{fig:amdahl}
\end{figure}

Chapters~2 and~3 have now made two design choices. Each core has a private cache
so that the first level stays fast (\cref{sec:cache-private}), and all cores share
one address space so that programs like \cref{lst:omp} are easy to write. Together
they allow one memory block to sit in several private caches at once, so the
system needs a rule for keeping the copies consistent (coherence) and hardware
that enforces the rule on every access (a coherence protocol). Chapter~4 covers
both.
